% =====================================================================
% III-C: cấp độ 3 (đánh giá lời giải thích + MMKG tiếng Việt). Kết quả ở Mục IV-D, IV-E.
% Nguồn số: experiments/xai-metrics-v2/{REPORT.md,summary.json}, experiments/xai-metrics-v3/REPORT.md,
% experiments/xai-ground-truth/REPORT.md; vietnamese-mmkg/DATA_CARD.md, experiments/checks/multilingual_clip_probe.json.
\subsection{Cấp độ 3: đánh giá lời giải thích và MMKG tiếng Việt}
\label{sec:method-l3}

Cấp độ 3 có hai phần. Phần thứ nhất đo fidelity (độ trung thực), validity (độ hợp lệ) và sparsity (độ gọn) của lời giải
thích ở Mục~\ref{sec:method-expl}, mở rộng phép thử xoá của Mục~\ref{sec:l1-cases} ra toàn bộ 10\,633 câu val, 3 seed mỗi
mô hình, cho cả mô hình chính \texttt{countonly\_factual\_r3} (gọi tắt là countonly) và \texttt{main\_r3} (cấp 1). Chỉ dùng
val: không train, không chấm test, không chỉnh mô hình theo kết quả đo. Ảnh được giải thích là ảnh mô hình xếp hạng 1 cho
mỗi câu, dù đúng hay sai. Phần thứ hai xây một MMKG tiếng Việt và chạy phương pháp tốt nhất của cấp~2 trên đó.

\subsubsection{Lời giải thích E và phép sửa đồ thị}\label{sec:xai-setup}
E được đọc ra từ chính lượt xếp hạng, không tốn thêm lượt chấm nào. Mỗi phần $j$ của ảnh hạng 1 (một node vật thể hoặc một
bộ ba có nhãn) nhận term
\begin{equation}
  t_j=\frac{\partial s}{\partial S_{\kappa}}\sum_i w_i\,a_{ij}\,\mathit{sim}_{ij},
  \label{eq:xai-term}
\end{equation}
với $\kappa\in\{\text{obj},\text{tri}\}$ là kênh của phần $j$, $S_\kappa$ là điểm kênh và $s$ là điểm cuối.
E là danh sách các phần xếp theo $t_j$ giảm dần; top-$k$ là $k$ phần đầu.
Với fusion raw của countonly, đạo hàm trong \eqref{eq:xai-term} là trọng số kênh ($b$ hoặc $c$ ở Mục~\ref{sec:method-main}),
nên các term cộng lại đúng bằng phần đồ thị của điểm cuối. Với fusion z-score của \texttt{main\_r3}, đạo hàm được tính chính
xác bằng autograd trên 50 ứng viên, để node và bộ ba được so trên cùng thang điểm cuối.

Mọi độ đo đều sửa đồ thị của ảnh hạng 1 rồi chấm lại bằng chính mô hình đó (\file{src/explain/xai_metrics.py}).
\emph{Bỏ một node} là bỏ node cùng mọi cạnh chạm vào nó, nên các bộ ba chạm vào nó cũng mất; cạnh kNN không được dựng lại.
\emph{Bỏ một bộ ba} là bỏ đúng cạnh có nhãn đó.
\emph{Chỉ giữ top-$k$} là giữ các phần đó cùng node tổng (nếu có); bộ ba được giữ thì hai node đầu--cuối của nó ở lại làm
``carrier'', vẫn truyền tin qua GAT nhưng không được chấm ở kênh vật thể. Ảnh sau khi sửa được mã hoá lại bằng cả mô hình,
kể cả GAT; 49 ứng viên còn lại giữ nguyên điểm kênh, còn fusion được tính lại trên cả 50 ứng viên.

Mốc \emph{gốc} là ảnh không sửa, chấm lại qua đúng đường chấm đó; ở cả 3 seed của cả hai mô hình, ảnh này vẫn hạng 1 ở mọi
câu, sai số điểm tối đa $2{,}4\cdot10^{-7}$ (countonly) và $2{,}0\cdot10^{-5}$ (\texttt{main\_r3}).
Mốc \emph{rỗng} bỏ mọi phần; \emph{hiệu ứng đồ thị} là $\Delta=s_{\text{gốc}}-s_{\text{rỗng}}$.
Mốc \emph{ngẫu nhiên} (NN) là trung bình 10 lần bốc cùng số node và cùng số bộ ba với top-$k$, rng cố định theo (seed, ảnh,
câu). Câu có E rỗng (31--36 câu mỗi seed với countonly, 39--45 với \texttt{main\_r3}) bị loại khỏi trung bình fidelity.
R@1 tính lại trong lúc đo trùng \file{metrics.json} của từng checkpoint (script kiểm bằng \texttt{assert}).

\subsubsection{Độ đo}\label{sec:xai-metrics}
\textbf{Fidelity} (tên theo~\cite{xai-yuan2023survey}), với $k\in\{1,3,5\}$.
Fid$^{+}$@$k$ bỏ top-$k$ của E và đếm tỉ lệ câu ảnh mất hạng 1; Fid$^{-}$@$k$ chỉ giữ top-$k$ và ghi hạng của ảnh.
Đường cong deletion / insertion~\cite{xai-petsiuk2018rise} bỏ dần, hoặc chỉ giữ dần, $k=1,\dots,10$ phần đầu;
$\mathrm{AUC}=\frac{1}{10}\sum_{k}(s_k-s_{\text{rỗng}})/\Delta$, deletion thấp là tốt, insertion cao là tốt.
AUC có thể âm hoặc lớn hơn 1, vì sửa đồ thị có thể đẩy điểm ra ngoài khoảng $[s_{\text{rỗng}},s_{\text{gốc}}]$.
Fid$_{\alpha}^{\pm}$~\cite{xai-zheng2024fidalpha} bỏ mỗi phần trong 5 phần đầu (Fid$_{\alpha}^{+}$), hoặc mỗi phần còn lại
(Fid$_{\alpha}^{-}$), với xác suất $\alpha=0{,}1$, rồi lấy phần $\Delta$ bị mất, trung bình 10 lần bốc. Đồ thị sau khi sửa
vẫn gần đồ thị gốc, và lời giải thích trung thực cho Fid$_{\alpha}^{+}>{}$Fid$_{\alpha}^{-}$.
\emph{Tính cộng} (additivity): vì E là một phép phân rã điểm, tổng term của top-$k$ là ``dự đoán'' của điểm mất khi bỏ
top-$k$; nhóm so nó với điểm mất đo được qua tỉ số tổng (tổng đo được chia tổng dự đoán, trên mọi câu) và tương quan
Spearman giữa các câu.

\textbf{Validity.}
Về phản thực tế (counterfactual), $k_{\text{flip}}$ là $k\le 10$ nhỏ nhất mà bỏ top-$k$ làm ảnh mất hạng 1.
Về ngữ nghĩa, với mỗi node vật thể của câu, nhóm xét node ảnh có $a_{ij}$ lớn nhất có cùng loại với nó không.
\emph{Strict}: cùng nhãn sau alias VG, hoặc synset của node ảnh trùng hay là hyponym của nghĩa danh từ đầu tiên của từ trong
câu; node người chỉ khớp với từ chung (person, people\ldots).
\emph{Lenient} là tiêu chí của bản đo v1: cùng nhãn sau alias, hoặc cùng loại theo \texttt{same\_type} (chấp nhận ``woman''
$\sim$ ``man'' và hypernym hai chiều), hoặc cosine CLIP $\ge 0{,}955$.
Dòng bộ ba khớp đủ khi chủ thể và đối tượng khớp strict và vị từ khớp sau alias.
Khớp nhóm số lượng chỉ chấm mô hình có bảng $g$.

\textbf{Sparsity}, theo phần của ảnh: $k_{80}$ là $k\le 10$ nhỏ nhất mà chỉ giữ top-$k$ còn ít nhất 80\% $\Delta$; $n_{80}$
là số phần đầu của E chiếm 80\% khối term dương.

Phép đo có ba bản: v1 (\file{experiments/xai-metrics/}), v2 (\file{experiments/xai-metrics-v2/}) là v1 sau khi soát lại
(sáu điểm sửa ghi ở Phụ lục~\ref{app:xai}), và v3 (\file{experiments/xai-metrics-v3/}) thêm ba cách giải thích. Mọi số ở phần kết quả là của v2
và v3, trừ các số v1 được nêu để đính chính; v3 giữ nguyên mọi số của v2 (14\,542 giá trị trùng tuyệt đối).

\subsubsection{Cách giải thích đối chứng và benchmark có đáp án biết trước}\label{sec:xai-baselines}
Nhóm đo thêm 15 cách giải thích khác trên cùng các phần của ảnh, bằng cùng các độ đo, trên một mẫu cố định 1\,000 câu val.
Bốn mốc đơn giản: ngẫu nhiên, E đảo ngược, cosine CLIP (không dùng mô hình) và gradient $\times$ input.
Hai cách thuộc tính chung: Integrated Gradients~\cite{xai-sundararajan2017ig} (IG, 20 bước) và
occlusion~\cite{xai-zeiler2014occlusion} (bỏ từng phần). Tám cách viết cho GNN: GNNExplainer~\cite{xai-ying2019gnnexplainer},
PGExplainer~\cite{xai-luo2020pgexplainer}, GNN-LRP~\cite{xai-schnake2022gnnlrp}, GraphSVX~\cite{xai-duval2021graphsvx},
SubgraphX~\cite{xai-yuan2021subgraphx} (v2), cùng phản thực tế CF-GNNExplainer~\cite{xai-lucic2022cfgnn}, cổng trên message
GraphMask~\cite{xai-schlichtkrull2021graphmask} và Shapley có cấu trúc GStarX~\cite{xai-zhang2022gstarx} (v3). Cuối cùng là
giá trị Shapley lấy mẫu 16 hoán vị~\cite{xai-castro2009shapley}, dùng làm chuẩn. Hàm giá trị của Shapley là biên giữa ảnh
được giải thích và ứng viên tốt nhất còn lại khi chỉ giữ một tập phần, theo ý tưởng của
RankingSHAP~\cite{xai-heuss2025rankingshap}; nó tốn khoảng 700--740 lượt chấm mỗi câu.

Các cách viết cho GNN vốn dành cho bài toán phân loại node hoặc đồ thị, nên được chỉnh cho các phần ảnh của reranker
(\file{src/explain/xai_baselines.py}). Kết quả của chúng vì vậy phụ thuộc cách chỉnh và ngân sách đã chọn, không phải kết
quả của bản cài gốc trên tác vụ gốc. Cách chỉnh từng cách và các phép kiểm tra cài đặt ở Phụ lục~\ref{app:xai}.

\textbf{Đáp án biết trước.} Validity ngữ nghĩa chấm bằng WordNet, không biết phần nào thực sự khiến ảnh đứng hạng 1. Nhóm
dựng ba bộ có đáp án sẵn, không gán tay, và chấm cả 16 cách bằng GEA~\cite{xai-agarwal2023graphxai}, Jaccard giữa $g$ phần
đầu của thứ tự và đáp án $g$ phần; mỗi mô hình, mỗi seed, 100 ca mà ảnh có đáp án đứng hạng 1.
(B) Dữ liệu tổng hợp: câu là một motif 2--3 vật, ảnh vàng chứa motif, và mỗi ảnh âm thiếu đúng một phần của motif, nên mô
hình cần mọi phần của đáp án~\cite{xai-faber2021groundtruth}; điểm CLIP là hằng số, và mô hình cùng kiến trúc countonly học
lại đạt R@1 100\% trên val tổng hợp.
(C) Cài một quan hệ của câu cùng một mồi nhử không liên quan vào một ứng viên hạng 2--5, ở vị trí ngẫu nhiên trong danh sách
phần của ảnh, chỉ giữ ca mà ảnh lên hạng 1.
(A) Mô tả vùng của Visual Genome làm truy vấn, đáp án là các object và quan hệ người gán nhãn nói tới; bộ này đo độ hợp lý,
không phải độ trung thực.

\subsubsection{MMKG tiếng Việt}\label{sec:vn-data}
Cấp độ 3 yêu cầu tự xây một MMKG nhỏ cho một miền hẹp, ưu tiên tiếng Việt, rồi chạy phương pháp tốt nhất của cấp 1--2 trên
đó. Data card (nguồn, cách tạo triple, giấy phép, điều kiện sử dụng) ở \file{vietnamese-mmkg/DATA_CARD.md}; thống kê đầy đủ,
mẫu dữ liệu và một ví dụ suy luận từng bước ở Phụ lục~\ref{app:vnmmkg}.

\textbf{Miền và nguồn.} Di tích và công trình ở Việt Nam (chùa, đền, tháp, nhà thờ, lăng, bảo tàng, nơi có danh hiệu di
sản), vì miền này có dữ liệu mở và nhiều quan hệ không thấy được trong ảnh: ở tỉnh nào, xây thời nào, ai sáng lập, thờ ai.
Thực thể và triple lấy từ Wikidata~\cite{vn-vrandecic2014wikidata} (CC0); phần mở đầu và infobox lấy từ Wikipedia tiếng
Việt~\cite{vn-viwiki} (CC BY-SA 4.0); ảnh lấy từ Wikimedia Commons~\cite{vn-commons}, tải ngày 24--25/10/2026.

\textbf{Quy mô} (khung gợi ý của đề trong ngoặc): 150 thực thể (50--150), gồm 75 di tích và 75 thực thể ngữ cảnh (tỉnh,
người, tôn giáo, danh hiệu, thời kỳ, lớp,~\ldots); 325 triple (300--1\,000) thuộc 13 quan hệ; 331 ảnh (200--500), mỗi di
tích 3--5 ảnh. Di tích là thực thể Wikidata ở Việt Nam có danh hiệu di sản hoặc thuộc một trong 10 lớp công trình, xếp theo
số sitelink.

\textbf{Triple.} Mỗi bộ ba giữ một nguồn, ưu tiên Wikidata, rồi infobox, rồi suy diễn. 257 triple là giá trị truthy (hạng
tốt nhất) trên Wikidata; 26 triple đọc từ ba trường infobox (người sáng lập, thờ, tôn giáo), chỉ nhận giá trị là wikilink
để ánh xạ được sang QID; 42 triple sinh từ hai quy tắc cố định: tỉnh (theo chuỗi \texttt{P131}, chỉ tạo khi tới đúng một
tỉnh) và thời kỳ (năm dựng đối chiếu với bảng 8 mốc triều đại). Toàn bộ 107 triple từ infobox và suy diễn được duyệt tay;
1 triple bị loại.

\textbf{Liên kết ảnh -- thực thể.} Ứng viên là ảnh \texttt{P18} cộng các file nằm trực tiếp trong Commons category của di
tích (JPEG/PNG rộng từ 400\,px, tối đa 12). Từng ảnh được duyệt tay, chỉ giữ ảnh thể hiện chính di tích; bản đồ, ảnh gần
trùng, ảnh có người là chủ thể chính bị bỏ (541 ảnh). Ảnh dùng chung chỉ giữ ở di tích xếp trên, nên mỗi ảnh thuộc đúng
một di tích; 35/110 di tích bị loại vì còn dưới 3 ảnh.

\textbf{Truy vấn.} 197 câu \emph{tự động} là tối đa 3 câu đầu bài Wikipedia, đã che nhãn và alias. 30 câu \emph{soạn tay}
gồm 12 câu trực tiếp (tả điều thấy trong ảnh) và 18 câu gián tiếp (cần ít nhất một quan hệ không thấy trong ảnh); đáp án
câu gián tiếp được tính tự động bằng cách khớp mẫu quan hệ trên triple. Câu đúng ở hạng $k$ nếu trong $k$ ảnh đầu có ảnh
của một di tích đáp án.

\textbf{Giấy phép và kiểm tra.} Chỉ nhận ảnh CC0, Public Domain, CC BY, CC BY-SA (163 ảnh CC BY-SA, 83 CC BY, 57 Public
Domain, 28 CC0); tác giả, giấy phép và URL nguồn từng ảnh ghi trong \file{images.jsonl}. Bộ dữ liệu chỉ dùng cho học thuật
trong môn học, phải ghi công tác giả khi hiển thị ảnh. Quyết định duyệt ảnh và triple lưu trong \file{review_images.json},
\file{review_triples.json}; chạy lại \texttt{python -m src.data.build\_vnmmkg} từ cache cho ra đúng các file đã commit.

\textbf{Giao thức đánh giá.}\label{sec:vn-leak} Bộ này quá nhỏ để vừa chọn tham số vừa chấm điểm, nên nhóm chốt hai giao
thức trước khi chạy.
\begin{itemize}
\item \textbf{A0, zero-shot (chính):} không chọn gì trên bộ tiếng Việt. Pool 331 ảnh, 197 câu tự động và 30 câu soạn tay.
  Encoder chọn trên MKG-W val (Mục~\ref{sec:vn-encoder}); $\alpha=0$ lấy từ \file{experiments/level2/val/choice.json} (hàng
  \texttt{kg\_triple}, cả ba seed).
\item \textbf{A1, chọn lại $\alpha$ (phụ):} chia di tích 50/50 bằng seed 0 (val: 38 di tích, 170 ảnh, 94 câu; test: 37 di
  tích, 161 ảnh, 103 câu). $\alpha$ quét từ 0 đến 1, bước 0{,}05, chọn theo R@1 trên val (hoà thì theo MRR); test chạy một
  lần. Câu soạn tay không tham gia.
\item Câu soạn tay chỉ dùng để chấm: \file{queries_hand.json} khoá ở commit \texttt{55fa697} (26/10/2026), trước mọi lần
  chạy, không sửa sau đó.
\item Triple là tri thức có sẵn, không phải cạnh thống kê đồng xuất hiện, nên được dùng hết.
\item Encoder đóng băng, không huấn luyện: kết quả tất định, chạy một lần, không có $\pm$std. Độ bất định đến từ cỡ mẫu:
  một câu bằng 5{,}56 điểm ở nhóm gián tiếp, 8{,}33 điểm ở nhóm trực tiếp, 0{,}97 điểm ở test A1.
\end{itemize}

\textbf{Phương pháp chuyển giao.}\label{sec:vn-method} Phương pháp được chuyển là kênh đếm cạnh KG (hàng \texttt{kg\_triple},
Mục~\ref{sec:method-l2}), mô hình tốt nhất ở cấp 2. Bộ xếp hạng GAT của cấp 1 cần scene graph của từng ảnh, còn bộ này chỉ có
KG mức thực thể. Kênh embedding NativE phải train lại trên 325 triple, quá ít để học embedding; hơn nữa trên MKG-W val kênh
này nhận trọng số $\beta=0$, nên mô hình chính của cấp 2 trùng với \texttt{kg\_triple}. Chỉ thay đồ thị, ngôn ngữ và encoder.
Linker (\texttt{EntityLinker}) chuyển câu về chữ thường, tách theo chữ cái (mỗi âm tiết là một token, giữ dấu), so khớp
chính xác với nhãn và alias tiếng Việt; độ dài tên tối đa nâng từ 5 token lên 12 âm tiết. Điểm đồ thị của một ảnh trong
top-50 là số cạnh KG trực tiếp giữa di tích của ảnh và các thực thể được nhắc; điểm cuối là
$\alpha\, z(\mathrm{CLIP}) + (1-\alpha)\, z(\text{đồ thị})$. Với $\alpha=0$, thứ tự CLIP chỉ còn dùng để phá thế hoà. Giải
thích là chính danh sách cạnh đã đếm. Code: \file{scripts/eval_vnmmkg_transfer.py}.

\subsubsection{Ablation}\label{sec:l3-ablation}
\textbf{Cạnh mà GAT chú ý nhất.} Đề yêu cầu bỏ hoặc nhiễu cả các cạnh dùng làm lời giải thích; ngoài các bộ ba (cạnh có
nhãn) trong E, nhóm thử cả các cạnh mà GAT chú ý nhất: với 3 node đầu của E, cắt đường truyền tin của các cạnh vào có
attention lớn nhất (giữ bộ ba), rồi so với cắt cùng số cạnh vào ngẫu nhiên của từng node. Trên 2\,803 câu mỗi seed có nhiều
cạnh vào hơn số cạnh cắt, tỉ lệ mất hạng 1 là \pmstd{1,0}{0,6}\% so với \pmstd{0,7}{0,4}\%, và điểm đổi
\pmstd{$-0{,}0011$}{0,0007} so với \pmstd{$-0{,}0010$}{0,0005} (\texttt{main\_r3}): cắt cạnh có attention cao không làm thay
đổi kết quả nhiều hơn cắt cạnh ngẫu nhiên. Với countonly (10\,481 câu mỗi seed), cắt cạnh có attention cao làm mất hạng 1 ở
\pmstd{3,5}{1,5}\% câu, khoảng gấp đôi cắt ngẫu nhiên (\pmstd{1,8}{0,9}\%) ở cả 3 seed, nhưng điểm đổi như nhau
(\pmstd{$+0{,}0025$}{0,0006} so với \pmstd{$+0{,}0021$}{0,0003}). Ở cả hai mô hình, hiệu ứng này nhỏ so với bỏ phần trong E;
trọng số attention của GAT vì vậy không xác định được cạnh quyết định thứ hạng, tương tự nhận xét
của~\cite{xai-jain2019attention}.

\textbf{Mức message.} Ngừng mọi message GAT của ảnh hạng 1 (giữ bộ ba) làm ảnh mất điểm ở 12,2\% câu với countonly, và
trung bình điểm còn tăng ($+0{,}0079$); với \texttt{main\_r3}, tỉ lệ là 48,2\% nhưng điểm chỉ giảm trung bình 0,0010. Ngừng
dần top-$k$ message theo GraphMask, theo attention GAT hay ngẫu nhiên: với countonly, không thứ tự nào hơn ngẫu nhiên; với
\texttt{main\_r3}, GraphMask và attention ngang nhau (AUC mất điểm 0,41 và 0,41) và hơn ngẫu nhiên (0,32). Truyền tin GAT
mang rất ít điểm, nên các cách giải thích theo message có ít thứ để giải thích trên mô hình này.

\textbf{Chọn encoder cho truy vấn tiếng Việt.}\label{sec:vn-encoder} CLIP ViT-B/32 của OpenAI (encoder cấp 1--2) chỉ học văn
bản tiếng Anh. Nhóm so bốn encoder trên 84 thực thể MKG-W val có mô tả tiếng Việt, mỗi thực thể hỏi bằng cả tiếng Việt và
tiếng Anh, pool 892 ảnh (Bảng~\ref{tab:vn-encoder}). Encoder đa ngữ gồm XLM-R~\cite{vn-conneau2020xlmr} ghép tháp ảnh huấn
luyện trên LAION-5B~\cite{vn-schuhmann2022laion5b} (bản OpenCLIP~\cite{vn-ilharco2021openclip}) và một bộ mã hoá văn bản đa
ngữ chưng cất vào không gian CLIP B/32~\cite{vn-reimers2020multilingual}. Với câu tiếng Việt, R@1 của OpenAI B/32 giảm từ
35,71 xuống 10,71, nên không dùng được. XLM-R large + ViT-H/14 cao nhất và không chênh giữa hai ngôn ngữ; nhóm chốt encoder
này. Hai encoder B/32 đa ngữ chỉ cách nhau 3 câu (mức nhiễu), còn H/14 hơn chúng khoảng 12 điểm R@1 (10 câu). Cấp 1--2 vẫn
giữ OpenAI B/32, nên cột ``CLIP'' ở kết quả tiếng Việt không so trực tiếp được với MKG-W.

\begin{table}[t]
\centering
\caption{Chọn encoder trên MKG-W val (84 thực thể, pool 892 ảnh). Tất định, chạy một lần. MRR nhân 100.}
\label{tab:vn-encoder}
\footnotesize
\setlength{\tabcolsep}{3pt}
\begin{tabular}{@{}lrrrrr@{}}
\toprule
 & \multicolumn{3}{c}{Tiếng Việt} & \multicolumn{2}{c}{Tiếng Anh} \\
\cmidrule(lr){2-4}\cmidrule(l){5-6}
Encoder (văn bản + ảnh) & R@1 & R@10 & MRR & R@1 & MRR \\
\midrule
OpenAI ViT-B/32 (cấp 1--2) & 10,71 & 17,86 & 14,10 & 35,71 & 46,03 \\
XLM-R base + ViT-B/32 & 25,00 & 61,90 & 37,83 & 28,57 & 40,50 \\
Chưng cất đa ngữ + ViT-B/32 & 28,57 & 55,95 & 40,05 & 33,33 & 44,62 \\
XLM-R large + ViT-H/14 & \textbf{40,48} & \textbf{78,57} & \textbf{51,80} & 40,48 & 51,31 \\
\bottomrule
\end{tabular}
\tabsrc{\file{experiments/checks/multilingual_clip_probe.json}}
\end{table}
